\chapter{Cayley's Theorem}
I am not happy with many textbook proofs of this theorem.
Some presentations introduce extra definitions too early
(for example, group actions), which can be an unnecessary burden
for beginners.

Many group theory books are very abstract. They introduce many
definitions quickly, and the motivation is not always clear.
Sometimes a new group is defined without clearly stating
what the multiplication is.

Think of me as a compiler: if a proof is not precise, it does not compile.
A programmer improves code until it compiles. In mathematics, readers are
often expected to adapt to unclear writing instead.

So I decided to prove the theorem myself until I fully understand it.


\section{The Theorem}
Every group is isomorphic to a collection of permutations.

\section{Proof}
We will construct a permutation from a given group. 
Every group has a multiplication table. We define permutations
from the columns of that table.

To keep notation simple, number the group's elements
\(1\) through \(n\), with \(1\) as the identity.
Also number the permutations from \(p_1\) through \(p_n\).

Next we create a multiplication table from the permutations.
\[
p_1=(1, 2, 3, 4 \dots n)
\]
being the identity permutation.
\[
p_i = (i, 2i, 3i, 4i, \dots ni)
\]
Consider the cell in the column \(p_i\) and row \(p_j\).
It is the permutation obtained by applying \(p_j\) to \(p_i\).
Let the result be \(p_k\). We will prove
\[k = i \cdot j\]

We call \(p_{i,m}\) the \(m\)-th element of the permutation \(p_i\).
The first element of \(p_k\) is the \(p_{j,1}\)-th element of \(p_i\).
The second element of \(p_k\) is the \(p_{j,2}\)-th element of \(p_i\).
...

The elements of \(p_j\) are \((j, 2j, 3j, 4j, \dots nj)\), meaning
\[
p_{j,m} = mj
\]
The elements of \(p_i\) are \((i, 2i, 3i, 4i, \dots ni)\).

\[
p_i\cdot p_j = (ij, 2ij, 3ij, 4ij, \dots nij)
\]

Therefore, any equation like \(p_i\cdot p_j = p_k\) mirrors
an equation in the original group, \(i\cdot j = k\).

\section{Numerical Example}
We start with a group given by the following table:

\begin{center}
    \begin{tabular}{|c|c|c|c|c|c|c|}
    \hline
     & \textbf{1} & \textbf{2} & \textbf{3} & \textbf{4} \\
    \hline
    \textbf{1} & 1 & 2 & 3 & 4 \\
    \hline
    \textbf{2} & 2 & 3 & 4 & 1 \\
    \hline
    \textbf{3} & 3 & 4 & 1 & 2 \\
    \hline
    \textbf{4} & 4 & 1 & 2 & 3 \\
    \hline
    \end{tabular}
\end{center}
We get the permutations from the columns:
\[
p_1=(1,2,3,4)\\
p_2=(2,3,4,1)\\
p_3=(3,4,1,2)\\
p_4=(4,1,2,3)\\
\]

Now check, for example, the result of \(p_2\cdot p_3\).
\[
p_2\cdot p_3 = (2,3,4,1) \cdot (3,4,1,2) = (4,1,2,3)
\]
In the original table we have:
\[
2 \cdot 3 = 4\\
2 \cdot 2 \cdot 3 = 1\\
3 \cdot 2 \cdot 3 = 2\\
4 \cdot 2 \cdot 3 = 3\\
\]

You can repeat the same check for all other cells.
So the original table can indeed be reconstructed from the permutations.
\begin{table}[ht]
\centering
    \begin{tabular}{|c|c|c|c|c|c|c|}
    \hline
     & \textbf{(\textbf{1},2,3,4)} & \textbf{(\textbf{2},3,4,1)} & \textbf{(\textbf{3},4,1,2)} & \textbf{(\textbf{4},1,2,3)} \\
    \hline
    \textbf{(\textbf{1},2,3,4)} & (\textbf{1},2,3,4) & (\textbf{2},3,4,1) & (\textbf{3},4,1,2) & (\textbf{4},1,2,3) \\
    \hline
    \textbf{(\textbf{2},3,4,1)} & (\textbf{2},3,4,1) & (\textbf{3},4,1,2) & (\textbf{4},1,2,3) & (\textbf{1},2,3,4) \\
    \hline
    \textbf{(\textbf{3},4,1,2)} & (\textbf{3},4,1,2) & (\textbf{4},1,2,3) & (\textbf{1},2,3,4) & (\textbf{2},3,4,1) \\
    \hline
    \textbf{(\textbf{4},1,2,3)} & (\textbf{4},1,2,3) & (\textbf{1},2,3,4) & (\textbf{2},3,4,1) & (\textbf{3},4,1,2) \\
    \hline
    \end{tabular}
    \caption{The multiplication table from the permutations. Using 
    the first element of each permutation, we recover the original multiplication table.}
\end{table}

